\documentclass[12pt]{article}
\usepackage{amsmath, amssymb, booktabs}
\usepackage[margin=1in]{geometry}

\begin{document}

\section{Parameter Estimates}
\input{parameter_estimates.tex}

\section{Moment Comparison}
\input{moments_comparison.tex}


\subsection*{Explanation of Parameter Estimates}
The table above shows the parameter estimates from the two-step SMM procedure. 
Both the identity and weighted estimates end up 
the same. This happens because the second-step weighting matrix does not change 
the shape of the objective function very much.

The estimates for $\alpha$, $\gamma$, $\rho$, and $\sigma$ are all positive and 
relatively small. Their t-statistics are quite large, meaning the model is very 
sensitive to small changes in these parameters around the optimum. The estimate 
for $\phi_0$ is exactly zero, which means the model does not need extra financial 
frictions to match the data under this specification.

\subsection*{Explanation of Moment Comparison}
The moment comparison table shows that the simulated moments are all equal to 
zero. This happens because, under the estimated parameters, the model simulation 
produces almost no variation in investment, profits, or $Q$. In other words, the 
policy functions lead to almost flat behavior, so the simulated panel of firms 
does not move much.

Since the simulated series are basically constant, their standard deviations and 
autocorrelations become zero by construction. Because of this, the model cannot 
match the empirical targets, and the differences are simply the negative of the 
targets. This suggests that the model might be too restrictive, or the SMM 
estimates push the model toward a corner solution with very little dynamic 
behavior.

\section*{Model Description}

The model follows the setup in Cooper and Ejarque (2003). Each firm chooses next
period's capital $K'$ to maximize the expected discounted value of future profits.
Output is produced using a Cobb--Douglas function with a productivity shock.
Productivity follows an AR(1) process in logs, which we approximate on a 
discrete grid using Tauchen's method.

Firms face convex adjustment costs when changing capital. These costs depend on
the size of investment relative to the capital stock. The parameter $\phi_0$
controls an external finance cost, but in our estimates it ends up being zero.
Given the model parameters, we solve the firm's dynamic programming problem by
value function iteration and obtain decision rules for capital.

Using these policy rules, we simulate a panel of firms and compute model moments
that correspond to the moments reported in Table 3 of Cooper and Ejarque.
We then estimate the parameters $(\alpha, \gamma, \rho, \sigma, \phi_0)$ using a 
two--step SMM procedure. The first step uses the identity weighting matrix, and 
the second step uses an estimated optimal weighting matrix based on the 
variance of the simulated moments.


\end{document}
